\documentclass[../DoAn.tex]{subfiles}
\begin{document}

Chương này trình bày khảo sát hiện trạng và phân tích yêu cầu cho hệ thống. Nội dung chương bám theo cấu trúc phân tích thiết kế hướng đối tượng của template SOICT: khảo sát hiện trạng, tổng quan chức năng, đặc tả chức năng và yêu cầu phi chức năng. Các yêu cầu được xây dựng từ bài toán gốc là trực quan hóa GNSS và khai thác Optical Flow trên thiết bị Android, sau đó mở rộng thành luồng xử lý video RAFT trên máy chủ.

\section{Khảo sát hiện trạng}
\label{section:2.1}
\subsection{Nhóm ứng dụng bản đồ và định vị}
\label{subsection:2.1.1}
Các ứng dụng bản đồ phổ thông như Google Maps, Apple Maps hoặc các ứng dụng dựa trên OpenStreetMap đã giải quyết tốt nhu cầu tìm kiếm địa điểm, hiển thị vị trí hiện tại và dẫn đường. Điểm mạnh của nhóm ứng dụng này là dữ liệu bản đồ phong phú, giao diện dễ dùng và khả năng tối ưu tuyến đường. Tuy nhiên, các ứng dụng này thường chỉ trình bày kết quả định vị cuối cùng mà không biểu diễn quá trình GNSS phía sau. Người dùng ít thấy được vệ tinh đang được quan sát, cường độ tín hiệu, chòm vệ tinh hoặc nguồn dữ liệu quỹ đạo. Vì vậy, nhóm ứng dụng này chưa đáp ứng tốt nhu cầu học tập và kiểm tra kỹ thuật liên quan đến GNSS.

Một số ứng dụng chuyên GNSS như GPSTest có thể hiển thị sky plot, danh sách vệ tinh, C/N0, azimuth và elevation. Các ứng dụng này phù hợp với việc kiểm tra tín hiệu vệ tinh nhưng thường tách khỏi trải nghiệm bản đồ hoặc không tích hợp với xử lý ảnh. Ngoài ra, nhiều ứng dụng chỉ hiển thị vệ tinh theo góc nhìn tương đối, chưa cung cấp một mô hình trực quan ba chiều hoặc AR kết hợp dữ liệu vị trí vệ tinh từ nhiều nguồn. Khoảng trống này là cơ sở để đồ án thiết kế module GNSS vừa có bản đồ, vừa có mô hình 3D/AR, vừa có thông tin nguồn dữ liệu.

\subsection{Nhóm ứng dụng xử lý ảnh và Optical Flow}
\label{subsection:2.1.2}
Trong lĩnh vực xử lý ảnh, Optical Flow thường được trình diễn bằng các ví dụ độc lập: đọc camera, tính vector và vẽ lên ảnh. OpenCV cung cấp các hàm như \texttt{goodFeaturesToTrack}, \texttt{calcOpticalFlowPyrLK} và \texttt{calcOpticalFlowFarneback}, giúp việc xây dựng demo trở nên thuận tiện \cite{opencv-video-tracking}. Tuy nhiên, các demo này thường thiếu giao diện điều khiển hoàn chỉnh, thiếu cơ chế ghi video, thiếu lưu phiên đo và thiếu so sánh giữa nhiều thuật toán trong cùng điều kiện.

KLT và Farneback là hai thuật toán có tính chất khác nhau. KLT phù hợp với môi trường có nhiều điểm đặc trưng rõ, trong khi Farneback có khả năng mô tả chuyển động dày đặc hơn trên toàn ảnh. Nếu chỉ chọn một thuật toán, hệ thống sẽ khó giúp người dùng hiểu các đánh đổi về hiệu năng và chất lượng hiển thị. Vì vậy, đồ án cần một cơ chế so sánh trực tiếp giữa hai thuật toán, kèm các chỉ số đo được trong quá trình chạy.

\subsection{Nhóm giải pháp xử lý video bằng trí tuệ nhân tạo}
\label{subsection:2.1.3}
Các mô hình học sâu cho Optical Flow như RAFT có khả năng biểu diễn trường chuyển động dày đặc và phức tạp hơn các thuật toán cổ điển \cite{raft2020}. Trong phạm vi đồ án, RAFT được đặt ở phía backend server; ứng dụng di động chỉ gửi video, theo dõi tiến độ và tải kết quả.

Với hướng tiếp cận client-server, vấn đề không chỉ nằm ở thuật toán mà còn ở trải nghiệm xử lý video. Nếu một video dài được gửi qua một request đồng bộ, người dùng có thể gặp timeout hoặc không biết tiến độ. Do đó, backend xử lý video cần cơ chế job bất đồng bộ, upload theo chunk cho file lớn, hủy job khi người dùng dừng thao tác và tải kết quả theo chunk khi file đầu ra lớn. Ứng dụng cần một worker chạy nền ở chế độ foreground, thông báo hệ thống và overlay để giữ thông tin tiến độ ngay cả khi người dùng rời màn hình hiện tại.

\subsection{So sánh và đánh giá tổng hợp}
\label{subsection:2.1.4}
Bảng~\ref{table:survey} tóm tắt các nhóm giải pháp liên quan và khoảng trống mà đồ án hướng tới. Qua khảo sát, có thể thấy chưa có nhiều ứng dụng kết hợp đồng thời GNSS, bản đồ 3D/AR, Optical Flow thời gian thực, analytics và xử lý video RAFT trong một hệ thống thử nghiệm thống nhất.

\begin{table}[H]
\centering
\begin{tabular}{|p{3.2cm}|p{5.0cm}|p{5.2cm}|}
\hline
\textbf{Nhóm giải pháp} & \textbf{Khả năng chính} & \textbf{Hạn chế so với mục tiêu đồ án} \\ \hline
Ứng dụng bản đồ phổ thông & Hiển thị vị trí, tìm kiếm địa điểm, dẫn đường và vẽ tuyến đường. & Ít thể hiện trạng thái vệ tinh, nguồn quỹ đạo và dữ liệu GNSS thô. \\ \hline
Ứng dụng GNSS chuyên dụng & Hiển thị chòm vệ tinh, SVID, C/N0, azimuth, elevation và trạng thái used-in-fix. & Thường không kết hợp với phân tích chuyển động camera hoặc xử lý Optical Flow. \\ \hline
Demo Optical Flow & Minh họa vector chuyển động từ camera hoặc video bằng thuật toán cổ điển. & Thiếu chức năng lưu phiên đo, so sánh KLT/Farneback và kết nối với luồng xử lý AI ngoại tuyến. \\ \hline
Máy chủ suy luận AI & Có thể chạy mô hình lớn như RAFT trên CPU/GPU mạnh hơn thiết bị di động. & Cần API quản lý job, theo dõi tiến độ, tải kết quả và đồng bộ với trải nghiệm Android. \\ \hline
\end{tabular}
\caption{Tổng hợp khảo sát các nhóm giải pháp liên quan}
\label{table:survey}
\end{table}

\section{Tổng quan chức năng}
\label{section:2.2}
\subsection{Biểu đồ use case tổng quát}
\label{subsection:2.2.1}
Các use case tổng quát được mô tả trong Bảng~\ref{table:usecase-overview} và minh họa bằng biểu đồ use case ở Hình~\ref{fig:usecase-tong-quat}. Các tác nhân chính của hệ thống gồm người dùng Android, thiết bị Android cung cấp dữ liệu GNSS, camera và IMU, máy chủ Optical Flow và các dịch vụ dữ liệu bên ngoài như IGS, CelesTrak, OSRM/Nominatim.

\begin{table}[H]
\centering
\begin{tabular}{|p{3cm}|p{4.8cm}|p{5.5cm}|}
\hline
\textbf{Tác nhân} & \textbf{Use case chính} & \textbf{Mục đích} \\ \hline
Người dùng & Xem bản đồ GNSS hai chiều & Quan sát vị trí hiện tại, tìm địa điểm và xem tuyến đường. \\ \hline
Người dùng & Xem vệ tinh trên mô hình 3D/AR & Quan sát vệ tinh, Trái Đất ngày/đêm, Mặt Trăng và Mặt Trời trong không gian trực quan hơn bản đồ phẳng. \\ \hline
Người dùng & Live routing có hỗ trợ Optical Flow & Theo dõi tuyến đường và dùng camera/IMU hỗ trợ tạm thời khi GNSS không đủ tin cậy. \\ \hline
Người dùng & Phân tích camera bằng Optical Flow & Chạy KLT hoặc Farneback trên luồng camera, hiển thị vector/heatmap và motion trail. \\ \hline
Người dùng & Chạy phiên phân tích so sánh & So sánh KLT và Farneback theo FPS, thời gian xử lý và số vector hoạt động. \\ \hline
Người dùng & Xử lý video bằng RAFT & Chọn My Server hoặc Outer Server, gửi video lên backend, theo dõi tiến độ và mở video kết quả. \\ \hline
Backend xử lý video & Thực thi job RAFT & Nhận video, chạy suy luận ở phía server, ghi video đầu ra và trả kết quả qua API. \\ \hline
Dịch vụ dữ liệu ngoài & Cung cấp bản đồ, tuyến đường và quỹ đạo & Cung cấp dữ liệu bổ sung cho bản đồ, GNSS và mô hình vệ tinh. \\ \hline
\end{tabular}
\caption{Use case tổng quát của hệ thống}
\label{table:usecase-overview}
\end{table}

\pumlfig{figures/usecase_tong_quat.png}{Biểu đồ use case tổng quát của hệ thống}{fig:usecase-tong-quat}

\subsection{Biểu đồ use case phân rã module Quan sát vệ tinh}
\label{subsection:2.2.2}
Module Quan sát vệ tinh (GNSS Viewer) mô tả nhóm thao tác liên quan đến bản đồ, vị trí hiện tại, tuyến đường và dữ liệu vệ tinh. Nhóm này được phân rã thành các use case cụ thể trong Bảng~\ref{table:usecase-gnss-viewer} và biểu đồ use case ở Hình~\ref{fig:usecase-quan-sat-ve-tinh}.

\begin{table}[H]
\centering
\begin{tabular}{|p{2.2cm}|p{4.2cm}|p{6.6cm}|}
\hline
\textbf{Mã use case} & \textbf{Tên use case} & \textbf{Mục đích/Kết quả} \\ \hline
UC-GNSS-01 & Xem vị trí hiện tại & Hiển thị marker vị trí người dùng trên bản đồ hai chiều. \\ \hline
UC-GNSS-02 & Tìm kiếm địa điểm & Người dùng nhập từ khóa và chọn một địa điểm từ kết quả tìm kiếm. \\ \hline
UC-GNSS-03 & Vẽ tuyến đường & Tạo tuyến đường từ vị trí hiện tại tới điểm đích đã chọn. \\ \hline
UC-GNSS-04 & Xem danh sách vệ tinh & Hiển thị SVID, chòm vệ tinh, C/N0, elevation, azimuth và trạng thái used-in-fix. \\ \hline
UC-GNSS-05 & Xem vệ tinh trên mô hình 3D & Dựng Trái Đất 3D và vị trí vệ tinh theo nguồn dữ liệu khả dụng. \\ \hline
UC-GNSS-06 & Xem vệ tinh bằng AR & Phủ thông tin vệ tinh lên luồng camera theo hướng nhìn của thiết bị. \\ \hline
UC-GNSS-07 & Cập nhật dữ liệu quỹ đạo & Tải hoặc làm mới dữ liệu IGS, CelesTrak/TLE và dữ liệu thời gian phục vụ phân giải vệ tinh. \\ \hline
UC-GNSS-08 & Xem chi tiết vệ tinh & Hiển thị nguồn vị trí, thông số tín hiệu và trạng thái của vệ tinh được chọn. \\ \hline
\end{tabular}
\caption{Phân rã use case module Quan sát vệ tinh}
\label{table:usecase-gnss-viewer}
\end{table}

\pumlfig{figures/usecase_quan_sat_ve_tinh.png}{Biểu đồ use case phân rã module Quan sát vệ tinh}{fig:usecase-quan-sat-ve-tinh}

\subsection{Biểu đồ use case phân rã module Chỉ đường trực tiếp}
\label{subsection:2.2.3}
Module Chỉ đường trực tiếp (LiveRouting) mô tả luồng dẫn đường có hỗ trợ camera và IMU khi GNSS suy giảm. Nhóm này được phân rã thành các use case cụ thể trong Bảng~\ref{table:usecase-live-routing} và biểu đồ use case ở Hình~\ref{fig:usecase-chi-duong-truc-tiep}.

\begin{table}[H]
\centering
\begin{tabular}{|p{2.2cm}|p{4.2cm}|p{6.6cm}|}
\hline
\textbf{Mã use case} & \textbf{Tên use case} & \textbf{Mục đích/Kết quả} \\ \hline
UC-LR-01 & Chọn điểm đến & Người dùng chọn địa điểm cần đi tới để chuẩn bị dẫn đường. \\ \hline
UC-LR-02 & Bắt đầu dẫn đường & Ứng dụng hiển thị tuyến đường và bắt đầu theo dõi vị trí theo thời gian thực. \\ \hline
UC-LR-03 & Đánh giá chất lượng GNSS & Kiểm tra số vệ tinh, độ tin cậy vị trí và tuổi dữ liệu GNSS. \\ \hline
UC-LR-04 & Bật camera assist & Kích hoạt camera khi GNSS không đủ tin cậy để hỗ trợ ước lượng chuyển động. \\ \hline
UC-LR-05 & Ước lượng chuyển động bằng Optical Flow/IMU & Tính tốc độ, thay đổi hướng và độ tin cậy từ Optical Flow kết hợp yaw rate của IMU. \\ \hline
UC-LR-06 & Map matching theo route & Ghép vị trí ước lượng về tuyến đường hiện có khi khoảng cách, hướng và độ tin cậy phù hợp. \\ \hline
UC-LR-07 & Cập nhật lại tuyến khi lệch đường & Gọi lại dịch vụ route khi quỹ đạo ước lượng lệch khỏi tuyến hiện tại. \\ \hline
UC-LR-08 & Hiển thị path GNSS và path assist & Vẽ riêng đoạn đường GNSS, đoạn hỗ trợ bằng camera và đoạn kiểm thử dropout để quan sát lại. \\ \hline
\end{tabular}
\caption{Phân rã use case module Chỉ đường trực tiếp}
\label{table:usecase-live-routing}
\end{table}

\pumlfig{figures/usecase_chi_duong_truc_tiep.png}{Biểu đồ use case phân rã module Chỉ đường trực tiếp}{fig:usecase-chi-duong-truc-tiep}

\subsection{Biểu đồ use case phân rã module Camera Optical Flow}
\label{subsection:2.2.4}
Module Camera Optical Flow mô tả nhóm thao tác phân tích chuyển động camera thời gian thực. Nhóm này được phân rã thành các use case cụ thể trong Bảng~\ref{table:usecase-camera-optical-flow} và biểu đồ use case ở Hình~\ref{fig:usecase-camera-optical-flow}.

\begin{table}[H]
\centering
\begin{tabular}{|p{2.2cm}|p{4.2cm}|p{6.6cm}|}
\hline
\textbf{Mã use case} & \textbf{Tên use case} & \textbf{Mục đích/Kết quả} \\ \hline
UC-OF-01 & Chọn thuật toán KLT/Farneback & Người dùng chọn thuật toán Optical Flow cần chạy trên luồng camera. \\ \hline
UC-OF-02 & Chọn chế độ hiển thị & Chuyển giữa vector, heatmap hoặc motion trail tùy nhu cầu quan sát. \\ \hline
UC-OF-03 & Điều chỉnh độ nhạy & Thay đổi tham số xử lý để cân bằng giữa số điểm/vector và tốc độ xử lý. \\ \hline
UC-OF-04 & Chọn ROI & Người dùng khoanh vùng quan tâm để chỉ xử lý một phần khung hình. \\ \hline
UC-OF-05 & Bám đối tượng trong ROI & Theo dõi vùng ROI khi đối tượng hoặc camera dịch chuyển nhẹ. \\ \hline
UC-OF-06 & Ghi video kết quả & Lưu video đã phủ vector hoặc heatmap vào bộ nhớ thiết bị. \\ \hline
UC-OF-07 & Chụp ảnh kết quả & Lưu một frame kết quả để phục vụ minh họa hoặc kiểm tra. \\ \hline
UC-OF-08 & Bắt đầu/dừng phiên analytics & Chạy phiên đo để thu FPS, thời gian xử lý và số điểm/vector hoạt động. \\ \hline
\end{tabular}
\caption{Phân rã use case module Camera Optical Flow}
\label{table:usecase-camera-optical-flow}
\end{table}

\pumlfig{figures/usecase_camera_optical_flow.png}{Biểu đồ use case phân rã module Camera Optical Flow}{fig:usecase-camera-optical-flow}

\subsection{Biểu đồ use case phân rã module Phân tích hiệu năng Optical Flow}
\label{subsection:2.2.5}
Module Phân tích hiệu năng Optical Flow (Analytics) mô tả nhóm chức năng lưu, mở lại và so sánh kết quả phân tích Optical Flow. Nhóm này được phân rã thành các use case cụ thể trong Bảng~\ref{table:usecase-analytics} và biểu đồ use case ở Hình~\ref{fig:usecase-phan-tich-hieu-nang}.

\begin{table}[H]
\centering
\begin{tabular}{|p{2.2cm}|p{4.2cm}|p{6.6cm}|}
\hline
\textbf{Mã use case} & \textbf{Tên use case} & \textbf{Mục đích/Kết quả} \\ \hline
UC-AN-01 & Lưu phiên phân tích & Ghi metrics của phiên so sánh KLT/Farneback thành dữ liệu JSON. \\ \hline
UC-AN-02 & Xem danh sách phiên & Hiển thị các phiên analytics đã lưu trên thiết bị. \\ \hline
UC-AN-03 & Xem biểu đồ FPS & Trực quan hóa FPS theo thời gian trong phiên phân tích. \\ \hline
UC-AN-04 & Xem biểu đồ thời gian xử lý & Trực quan hóa process time của từng thuật toán. \\ \hline
UC-AN-05 & Xem số điểm/vector hoạt động & So sánh số track KLT hoặc vector Farneback theo từng sample. \\ \hline
UC-AN-06 & Xóa phiên phân tích & Xóa phiên đo không còn cần sử dụng khỏi bộ nhớ ứng dụng. \\ \hline
\end{tabular}
\caption{Phân rã use case module Phân tích hiệu năng Optical Flow}
\label{table:usecase-analytics}
\end{table}

\pumlfig{figures/usecase_phan_tich_hieu_nang.png}{Biểu đồ use case phân rã module Phân tích hiệu năng Optical Flow}{fig:usecase-phan-tich-hieu-nang}

\subsection{Biểu đồ use case phân rã module Xử lý video Optical Flow bằng AI}
\label{subsection:2.2.6}
Module Xử lý video Optical Flow bằng AI (Video RAFT Backend) mô tả luồng xử lý video ngoại tuyến bằng máy chủ Python tự triển khai hoặc backend ngoài. Nhóm này được phân rã thành các use case cụ thể trong Bảng~\ref{table:usecase-video-raft} và biểu đồ use case ở Hình~\ref{fig:usecase-xu-ly-video-ai}.

\begin{table}[H]
\centering
\begin{tabular}{|p{2.2cm}|p{4.2cm}|p{6.6cm}|}
\hline
\textbf{Mã use case} & \textbf{Tên use case} & \textbf{Mục đích/Kết quả} \\ \hline
UC-RFT-01 & Chọn video đầu vào & Người dùng chọn video cần xử lý từ bộ nhớ thiết bị. \\ \hline
UC-RFT-02 & Chọn backend và chế độ xử lý & Chọn My Server hoặc Outer Server, đồng thời chọn vector, heatmap hoặc ROI trước khi gửi video. \\ \hline
UC-RFT-03 & Upload video trực tiếp & Gửi video nhỏ bằng multipart request tới máy chủ. \\ \hline
UC-RFT-04 & Upload video theo chunk & Chia video lớn thành nhiều phần để tránh timeout và lỗi bộ nhớ. \\ \hline
UC-RFT-05 & Tạo job xử lý & Backend tạo \texttt{job\_id} hoặc mã theo dõi tương ứng và đưa video vào hàng xử lý nền. \\ \hline
UC-RFT-06 & Theo dõi tiến độ job & Ứng dụng poll trạng thái và hiển thị phần trăm xử lý cho người dùng. \\ \hline
UC-RFT-07 & Hủy job & Người dùng dừng job đang chạy khi không còn nhu cầu xử lý. \\ \hline
UC-RFT-08 & Tải video kết quả & Tải file đầu ra hoặc tải theo chunk khi video kết quả lớn. \\ \hline
UC-RFT-09 & Mở video đã xử lý & Mở video kết quả trong trình phát để kiểm tra vector hoặc heatmap. \\ \hline
UC-RFT-10 & Cleanup tài nguyên server & Gọi endpoint hoặc thao tác dọn file tạm sau khi tải xong, job lỗi hoặc job bị hủy. \\ \hline
\end{tabular}
\caption{Phân rã use case module Xử lý video Optical Flow bằng AI}
\label{table:usecase-video-raft}
\end{table}

\pumlfig{figures/usecase_xu_ly_video_ai.png}{Biểu đồ use case phân rã module Xử lý video Optical Flow bằng AI}{fig:usecase-xu-ly-video-ai}

\section{Đặc tả chức năng chi tiết}
\label{section:2.3}
Phần này lựa chọn năm use case quan trọng nhất của đồ án, tương ứng với năm module đã phân rã ở Mục~\ref{section:2.2}, để đặc tả chi tiết: Quan sát vệ tinh, Chỉ đường trực tiếp, Camera Optical Flow, Phân tích hiệu năng Optical Flow và Xử lý video Optical Flow bằng AI. Mỗi đặc tả trình bày các thông tin: tên use case, tác nhân, mô tả, tiền điều kiện, luồng sự kiện chính, luồng sự kiện phát sinh và hậu điều kiện. Với những use case có nhiều rẽ nhánh, luồng chính được trình bày kèm các điểm rẽ nhánh và bổ sung một biểu đồ hoạt động (activity) để minh họa trực quan luồng xử lý giữa người dùng và hệ thống.

\subsection{Đặc tả chức năng Quan sát vệ tinh}
\label{subsection:2.3.1}
\textbf{Tên use case:} Quan sát vệ tinh GNSS trên bản đồ, mô hình 3D và AR.

\textbf{Tác nhân:} Người dùng (tác nhân chính); thiết bị Android cung cấp vị trí và trạng thái GNSS, các dịch vụ dữ liệu ngoài Nominatim/OSRM cho bản đồ, IGS và CelesTrak cho dữ liệu quỹ đạo (tác nhân phụ).

\textbf{Mô tả:} Use case cho phép người dùng quan sát vị trí hiện tại trên bản đồ hai chiều, tìm kiếm địa điểm, vẽ tuyến đường, đồng thời quan sát vệ tinh trên mô hình Trái Đất ba chiều và màn hình thực tế tăng cường. Mỗi vệ tinh được hiển thị kèm nguồn phân giải vị trí để người dùng biết được độ tin cậy của dữ liệu.

\textbf{Tiền điều kiện:}
\begin{itemize}
    \item Người dùng đã cấp quyền vị trí hoặc đang chạy ở chế độ vị trí giả lập.
    \item GPS được bật; chế độ ba chiều yêu cầu thiết bị hỗ trợ OpenGL ES 3.2; chế độ AR yêu cầu ARCore.
\end{itemize}

\textbf{Luồng sự kiện chính:}
\begin{enumerate}
    \item Người dùng mở màn hình Quan sát vệ tinh.
    \item Hệ thống khởi tạo bản đồ osmdroid và đăng ký \texttt{LocationManager}, \texttt{GnssStatus.Callback} cùng \texttt{GnssMeasurementsEvent.Callback} nếu thiết bị hỗ trợ.
    \item Hệ thống hiển thị marker vị trí người dùng trên bản đồ.
    \item Mỗi khi nhận \texttt{GnssStatus}, \texttt{GnssSatelliteTracker} phân giải vị trí từng vệ tinh theo chuỗi ưu tiên: PVT thật từ thiết bị, IGS Broadcast Ephemeris, CelesTrak TLE/SGP4 và cuối cùng là phép xấp xỉ từ azimuth/elevation.
    \item Tại điểm rẽ nhánh này, người dùng có thể chọn một trong các thao tác:
    \begin{enumerate}[label=(\alph*)]
        \item Nếu chọn tìm kiếm địa điểm, chuyển sang bước 6.
        \item Nếu chọn xem mô hình 3D, chuyển sang bước 7.
        \item Nếu chọn xem AR, chuyển sang bước 8.
        \item Nếu chọn xem chi tiết một vệ tinh, chuyển sang bước 9.
    \end{enumerate}
    \item Người dùng nhập từ khóa; hệ thống gọi Nominatim, hiển thị kết quả và vẽ tuyến đường tới điểm đích bằng hình học OSRM. Quay lại bước 5.
    \item Hệ thống chuyển sang \texttt{EarthRenderer}, dựng Trái Đất ngày/đêm, Mặt Trời, Mặt Trăng và đặt vệ tinh theo nguồn dữ liệu khả dụng. Quay lại bước 5.
    \item Hệ thống mở màn hình AR và phủ thông tin vệ tinh lên luồng camera theo hướng nhìn của thiết bị. Quay lại bước 5.
    \item Hệ thống hiển thị chi tiết vệ tinh được chọn: chòm vệ tinh, SVID, C/N0, azimuth, elevation, trạng thái used-in-fix và nguồn phân giải vị trí.
\end{enumerate}

\textbf{Luồng sự kiện phát sinh:}

\textbf{Luồng 4a: Không trích xuất được PVT thật.}
\begin{enumerate}
    \item Hệ thống không lấy được PVT từ thiết bị.
    \item Hệ thống tự động lùi sang IGS, CelesTrak hoặc xấp xỉ azimuth/elevation và đánh dấu nguồn dữ liệu tương ứng.
\end{enumerate}

\textbf{Luồng 7a: Thiết bị không hỗ trợ OpenGL ES 3.2.}
\begin{enumerate}
    \item Hệ thống phát hiện thiết bị không đủ điều kiện dựng mô hình 3D.
    \item Hệ thống thông báo và giữ nguyên bản đồ hai chiều.
\end{enumerate}

\textbf{Luồng 4b: Không có vệ tinh khả dụng.}
\begin{enumerate}
    \item Sau một khoảng thời gian vẫn không nhận được vệ tinh nào.
    \item Hệ thống hiển thị hộp thoại cảnh báo thay vì dừng đột ngột.
\end{enumerate}

\textbf{Hậu điều kiện:}
\begin{itemize}
    \item Người dùng quan sát được vệ tinh kèm nguồn phân giải vị trí.
    \item Tuyến đường vừa tạo có thể được chuyển tiếp sang chức năng Chỉ đường trực tiếp.
\end{itemize}

\activitydiagram{figures/activity_quan_sat_ve_tinh.png}{Biểu đồ hoạt động chức năng Quan sát vệ tinh}{fig:activity-quan-sat-ve-tinh}

\subsection{Đặc tả chức năng Chỉ đường trực tiếp}
\label{subsection:2.3.2}
\textbf{Tên use case:} Dẫn đường theo tuyến có hỗ trợ Optical Flow và IMU khi GNSS suy giảm.

\textbf{Tác nhân:} Người dùng (tác nhân chính); thiết bị Android cung cấp GNSS, camera và IMU, dịch vụ định tuyến OSRM (tác nhân phụ).

\textbf{Mô tả:} Use case cho phép người dùng dẫn đường theo tuyến; khi GNSS suy giảm, hệ thống dùng camera kết hợp IMU để ước lượng chuyển động (visual odometry thực dụng) và map matching để giữ marker bám tuyến đường, đồng thời phân biệt rõ đoạn đường đi bằng GNSS và đoạn được hỗ trợ bằng camera.

\textbf{Tiền điều kiện:}
\begin{itemize}
    \item Đã có tuyến đường từ một điểm đến được chọn.
    \item Ứng dụng có quyền vị trí; có quyền camera nếu muốn bật hỗ trợ thị giác.
\end{itemize}

\textbf{Luồng sự kiện chính:}
\begin{enumerate}
    \item Người dùng chọn điểm đến và bắt đầu dẫn đường.
    \item Hệ thống tải tuyến đường, cập nhật marker theo GNSS và vẽ đoạn đường đã đi bằng GNSS (path màu đen).
    \item Hệ thống liên tục đánh giá chất lượng GNSS qua số vệ tinh used-in-fix, tuổi bản tin và độ chính xác của vị trí.
    \item Tại điểm rẽ nhánh này, hệ thống xét chất lượng GNSS:
    \begin{enumerate}[label=(\alph*)]
        \item Nếu GNSS còn tin cậy, hệ thống tiếp tục bám GNSS và hiệu chỉnh hệ số pixel/giây sang mét/giây cho Optical Flow; quay lại bước 3.
        \item Nếu GNSS không đủ tin cậy, chuyển sang bước 5.
    \end{enumerate}
    \item Hệ thống bật camera assist và chuyển sang chế độ ước lượng bằng cảm biến.
    \item Hệ thống tính tốc độ từ Optical Flow, kết hợp yaw rate của IMU cho hướng và cơ chế ZUPT để khử trôi khi dừng.
    \item Vị trí ước lượng được ghép về tuyến đường bằng map matching có điều kiện (route-lock); hệ thống vẽ đoạn hỗ trợ (path màu đỏ) và đánh dấu các điểm mất, khôi phục GNSS.
    \item Khi GNSS phục hồi, hệ thống tắt camera assist, nối lại đoạn đường GNSS và quay lại bước 3.
\end{enumerate}

\textbf{Luồng sự kiện phát sinh:}

\textbf{Luồng 5a: Người dùng từ chối quyền camera.}
\begin{enumerate}
    \item Ứng dụng vẫn dẫn đường bằng GNSS nhưng không bật camera assist.
\end{enumerate}

\textbf{Luồng 7a: Quỹ đạo ước lượng lệch tuyến qua nhiều mẫu.}
\begin{enumerate}
    \item Nếu có Internet, hệ thống gọi lại OSRM để cập nhật tuyến từ vị trí ước lượng tới đích.
    \item Nếu mất Internet, hệ thống giữ tuyến hiện tại và hiển thị thông báo offline.
\end{enumerate}

\textbf{Luồng 3a: Bật chế độ kiểm thử dropout.}
\begin{enumerate}
    \item Người dùng bật mô phỏng mất GNSS để quan sát lại hành vi hỗ trợ.
\end{enumerate}

\textbf{Hậu điều kiện:}
\begin{itemize}
    \item Marker được duy trì liên tục kể cả khi GNSS gián đoạn.
    \item Đoạn đường GNSS và đoạn hỗ trợ được phân biệt rõ trên bản đồ; phiên dẫn đường có thể được lưu lại để xem lại.
\end{itemize}

\activitydiagram{figures/activity_chi_duong_truc_tiep.png}{Biểu đồ hoạt động chức năng Chỉ đường trực tiếp}{fig:activity-chi-duong-truc-tiep}

\subsection{Đặc tả chức năng Camera Optical Flow}
\label{subsection:2.3.3}
\textbf{Tên use case:} Phân tích chuyển động camera thời gian thực bằng KLT hoặc Farneback.

\textbf{Tác nhân:} Người dùng (tác nhân chính); thiết bị Android cung cấp camera và IMU (tác nhân phụ).

\textbf{Mô tả:} Use case cho phép người dùng chạy KLT hoặc Farneback trên luồng camera thời gian thực, chọn chế độ hiển thị, điều chỉnh độ nhạy, khoanh đối tượng cần theo dõi (object tracker) và ghi lại kết quả chuyển động.

\textbf{Tiền điều kiện:}
\begin{itemize}
    \item Ứng dụng được cấp quyền camera.
    \item Thư viện OpenCV được nạp thành công.
\end{itemize}

\textbf{Luồng sự kiện chính:}
\begin{enumerate}
    \item Người dùng mở màn hình Camera Optical Flow và chọn thuật toán KLT hoặc Farneback.
    \item Người dùng chọn chế độ hiển thị (vector, heatmap hoặc motion trail) và điều chỉnh độ nhạy để cân bằng giữa số điểm/vector và tốc độ xử lý.
    \item CameraX gửi từng frame qua \texttt{ImageAnalysis} với chiến lược \texttt{STRATEGY\_KEEP\_ONLY\_LATEST}; hệ thống chuyển frame sang \texttt{Mat} của OpenCV, chạy thuật toán và phủ kết quả lên khung hình.
    \item Tại điểm rẽ nhánh này, người dùng có thể chọn:
    \begin{enumerate}[label=(\alph*)]
        \item Nếu khoanh một vùng quan tâm (ROI), chuyển sang bước 5.
        \item Nếu ghi video hoặc chụp ảnh, chuyển sang bước 6.
        \item Nếu không, tiếp tục quan sát và quay lại bước 3.
    \end{enumerate}
    \item Bộ object tracker (CSRT kết hợp template matching) bám vùng ROI khi camera hoặc đối tượng dịch chuyển; Optical Flow chỉ chạy trong vùng con. Quay lại bước 3.
    \item Hệ thống ghi video hoặc lưu ảnh kết quả vào bộ nhớ thiết bị.
\end{enumerate}

\textbf{Luồng sự kiện phát sinh:}

\textbf{Luồng 1a: Chưa cấp quyền camera.}
\begin{enumerate}
    \item Ứng dụng yêu cầu cấp quyền camera trước khi mở luồng.
\end{enumerate}

\textbf{Luồng 5a: Vùng ROI quá nhỏ.}
\begin{enumerate}
    \item Hệ thống yêu cầu chọn lại vì không đủ pixel để phát hiện đặc trưng hoặc lấy mẫu vector.
\end{enumerate}

\textbf{Luồng 3a: Chế độ tự động theo IMU.}
\begin{enumerate}
    \item Hệ thống dùng IMU để nhận biết trạng thái tĩnh hay di chuyển nhằm điều chỉnh cách hiển thị.
\end{enumerate}

\textbf{Hậu điều kiện:}
\begin{itemize}
    \item Kết quả Optical Flow được hiển thị thời gian thực.
    \item Video hoặc ảnh kết quả được lưu lại khi người dùng yêu cầu.
\end{itemize}

\activitydiagram{figures/activity_camera_optical_flow.png}{Biểu đồ hoạt động chức năng Camera Optical Flow}{fig:activity-camera-optical-flow}

\subsection{Đặc tả chức năng Phân tích hiệu năng Optical Flow}
\label{subsection:2.3.4}
\textbf{Tên use case:} So sánh hiệu năng KLT và Farneback trên cùng luồng camera.

\textbf{Tác nhân:} Người dùng (tác nhân chính); thiết bị Android cung cấp camera (tác nhân phụ).

\textbf{Mô tả:} Use case chạy đồng thời KLT và Farneback trên cùng một frame đầu vào để so sánh hiệu năng hai thuật toán và lưu lại phiên đo dưới dạng JSON, từ đó người dùng có thể mở lại và xem các biểu đồ.

\textbf{Tiền điều kiện:}
\begin{itemize}
    \item Camera đang hoạt động.
    \item Người dùng không ghi video song song tại thời điểm bắt đầu phân tích.
\end{itemize}

\textbf{Luồng sự kiện chính:}
\begin{enumerate}
    \item Người dùng nhấn bắt đầu phân tích.
    \item Hệ thống khởi tạo hai bộ thuật toán mới và khóa độ nhạy của cả hai ở mức tối đa để bảo đảm so sánh công bằng.
    \item Mỗi khung hình được nhân đôi và đưa vào hai thuật toán; kết quả được ghép trái--phải trên cùng màn hình để quan sát trực quan.
    \item Hệ thống ghi sample định kỳ (khoảng một mẫu mỗi 250 ms) gồm FPS, thời gian xử lý, số điểm/vector hoạt động, vector trung bình, độ tin cậy và ngưỡng của mỗi thuật toán.
    \item Người dùng dừng phân tích; hệ thống lưu \texttt{AnalyticsSession} thành tệp JSON trong bộ nhớ trong (\texttt{filesDir/analytics}).
    \item Người dùng mở lại phiên để xem các biểu đồ FPS, thời gian xử lý và số điểm/vector theo thời gian.
\end{enumerate}

\textbf{Luồng sự kiện phát sinh:}

\textbf{Luồng 1a: Đang ghi video.}
\begin{enumerate}
    \item Hệ thống yêu cầu dừng ghi video trước khi bắt đầu phân tích.
\end{enumerate}

\textbf{Luồng 5a: Phiên không có sample hợp lệ.}
\begin{enumerate}
    \item Hệ thống không tạo phiên rỗng.
\end{enumerate}

\textbf{Hậu điều kiện:}
\begin{itemize}
    \item Phiên đo được lưu dưới dạng JSON.
    \item Phiên có thể mở lại ở màn hình phân tích để xem biểu đồ.
\end{itemize}

\activitydiagram{figures/activity_phan_tich_hieu_nang.png}{Biểu đồ hoạt động chức năng Phân tích hiệu năng Optical Flow}{fig:activity-phan-tich-hieu-nang}

\subsection{Đặc tả chức năng Xử lý video Optical Flow bằng AI}
\label{subsection:2.3.5}
\textbf{Tên use case:} Xử lý video bằng mô hình RAFT trên backend server.

\textbf{Tác nhân:} Người dùng (tác nhân chính); máy chủ tự triển khai (My Server) và backend ngoài (Outer Server) (tác nhân phụ).

\textbf{Mô tả:} Use case xử lý các video đã chọn bằng mô hình RAFT trên máy chủ. Người dùng chọn backend và chế độ xử lý, gửi video, theo dõi tiến độ và mở video kết quả; tác vụ nền được quản lý độc lập với màn hình hiện tại để không gián đoạn khi người dùng chuyển màn hình.

\textbf{Tiền điều kiện:}
\begin{itemize}
    \item Cấu hình backend HTTPS đã được thiết lập.
    \item Thiết bị có kết nối mạng.
\end{itemize}

\textbf{Luồng sự kiện chính:}
\begin{enumerate}
    \item Người dùng chọn video đầu vào.
    \item Người dùng chọn backend (My Server hoặc Outer Server) và chế độ xử lý (vector, heatmap hoặc ROI).
    \item Tại điểm rẽ nhánh theo kích thước video:
    \begin{enumerate}[label=(\alph*)]
        \item Nếu video nhỏ, hệ thống gửi bằng multipart request; chuyển sang bước 5.
        \item Nếu video lớn, hệ thống tạo upload session, gửi từng chunk (kèm kiểm tra SHA-256) và ráp lại ở máy chủ.
    \end{enumerate}
    \item Máy chủ hoàn tất upload và tạo job.
    \item Máy chủ xử lý nền: chạy RAFT bằng ONNX Runtime, ghi MP4 H.264 bằng ffmpeg, và dùng Cutie để lan truyền mask đối tượng nếu có ROI.
    \item Ứng dụng theo dõi tiến độ qua một foreground worker (WorkManager) kèm thông báo hệ thống.
    \item Ứng dụng tải kết quả theo file hoặc theo chunk, mở bằng \texttt{VideoOpticalFlowFragment} và gọi dọn tài nguyên trên máy chủ.
\end{enumerate}

\textbf{Luồng sự kiện phát sinh:}

\textbf{Luồng 5a: Backend lỗi hoặc chưa nạp được mô hình.}
\begin{enumerate}
    \item Job chuyển sang trạng thái thất bại; ứng dụng hiển thị thông báo và không tạo file kết quả sai.
\end{enumerate}

\textbf{Luồng 6a: Người dùng rời màn hình.}
\begin{enumerate}
    \item Worker nền tiếp tục quản lý tiến độ; overlay giữ thông tin job để người dùng quay lại.
\end{enumerate}

\textbf{Luồng 6b: Người dùng hủy job.}
\begin{enumerate}
    \item Ứng dụng hủy phiên upload đang dở hoặc gửi yêu cầu hủy job tới máy chủ.
\end{enumerate}

\textbf{Luồng 4a: Hàng đợi máy chủ đầy.}
\begin{enumerate}
    \item Máy chủ trả về mã lỗi 429; ứng dụng thông báo và đề nghị thử lại sau.
\end{enumerate}

\textbf{Hậu điều kiện:}
\begin{itemize}
    \item Video kết quả được lưu vào thư viện của ứng dụng.
    \item Trường hợp lỗi, thông báo được trả về rõ ràng và tài nguyên tạm trên máy chủ được dọn.
\end{itemize}

\activitydiagram{figures/activity_xu_ly_video_ai.png}{Biểu đồ hoạt động chức năng Xử lý video Optical Flow bằng AI}{fig:activity-xu-ly-video-ai}

\section{Yêu cầu phi chức năng}
\label{section:2.4}
\subsection{Hiệu năng xử lý}
\label{subsection:2.4.1}
Về hiệu năng, các chức năng camera thời gian thực cần sử dụng chiến lược backpressure chỉ giữ frame mới nhất để tránh hàng đợi xử lý tăng không kiểm soát. KLT và Farneback cần có tham số độ nhạy để cân bằng giữa số điểm/vector và thời gian xử lý. Với RAFT, yêu cầu thời gian thực không đặt trên thiết bị Android; thay vào đó hệ thống phải có tiến độ job, giới hạn hàng đợi server, upload/download theo chunk khi cần và không giữ một HTTP request dài trong toàn bộ thời gian xử lý video.

\subsection{Độ ổn định và độ tin cậy}
\label{subsection:2.4.2}
Về độ tin cậy, ứng dụng cần xử lý các trường hợp thiếu quyền vị trí, GPS tắt, không có vệ tinh, không nạp được mô hình RAFT hoặc máy chủ trả lỗi. Các thao tác xử lý video phải tránh mất file đầu ra khi ứng dụng chuyển màn hình, vì vậy tiến độ được điều phối qua WorkManager foreground worker, bus sự kiện, notification và overlay độc lập với fragment hiện tại. Với GNSS, hệ thống cần có fallback khi PVT thật không tồn tại trên thiết bị.

\subsection{Tính dễ sử dụng}
\label{subsection:2.4.3}
Về tính dễ dùng, giao diện cần tách rõ các nhóm chức năng GNSS, Optical Flow, LiveRouting, Media và Analytics để người dùng dễ quan sát luồng xử lý. Các thao tác cần phản hồi trực tiếp bằng trạng thái tải, tiến độ xử lý, thông báo lỗi và kết quả hiển thị trên bản đồ, camera hoặc trình phát video. Với các quyền hệ thống như vị trí, camera và thông báo, ứng dụng cần hướng dẫn người dùng bật đúng quyền trước khi tiếp tục.

\subsection{Khả năng mở rộng và bảo trì}
\label{subsection:2.4.4}
Về khả năng bảo trì, mã nguồn cần được chia theo package chức năng như \texttt{gnss}, \texttt{optical\_flow}, \texttt{screen}, \texttt{video}, \texttt{model} và \texttt{util}. Các lớp thuật toán cần cùng thực thi interface \texttt{OpticalFlow} để có thể thay thế hoặc mở rộng. Máy chủ cần tách API ở \texttt{src/main.py} và xử lý suy luận ở \texttt{src/inference.py} để giảm phụ thuộc giữa giao thức HTTP và thuật toán.

\subsection{Bảo mật dữ liệu và quyền truy cập thiết bị}
\label{subsection:2.4.5}
Về bảo mật, ứng dụng yêu cầu URL máy chủ bắt đầu bằng HTTPS trong cấu hình chính thức. Các file video tạm trên máy chủ phải được dọn sau khi trả kết quả, khi client gọi cleanup hoặc khi job lỗi/hủy. Ứng dụng chỉ yêu cầu các quyền cần thiết cho camera, vị trí, mạng, foreground service và thông báo, đồng thời cần thông báo rõ ràng khi người dùng chưa cấp quyền hoặc chưa bật GPS.

\subsection{Tính tương thích trên thiết bị Android}
\label{subsection:2.4.6}
Về tính tương thích, ứng dụng cần xử lý khác biệt giữa các thiết bị Android về chất lượng GNSS, camera, cảm biến IMU, hiệu năng CPU/GPU và khả năng hỗ trợ thư viện. Những chức năng phụ thuộc phần cứng như dữ liệu GNSS thô hoặc AR cần có phương án fallback hoặc thông báo phù hợp khi thiết bị không đáp ứng điều kiện. Với video RAFT, ứng dụng chỉ cần cấu hình backend phù hợp và kết nối mạng ổn định.

\end{document}
